\chapter{Stanislav Smirnov --- Fields Medal 2010}\label{ch2010stanislavsmirnov:stanislav-smirnov-fields-medal-2010-a-comprehensive-report-on-his-work-and-impact}

\begin{center}\rule{0.5\linewidth}{0.5pt}\end{center}

\textbf{Prepared:} July 2026 \textbf{Subject:} Stanislav Smirnov (b.
September 3, 1970, Leningrad, USSR) \textbf{Position:} Professor,
University of Geneva; Head of the Laboratory of Mathematical Physics,
Chebyshev Laboratory, St.~Petersburg State University \textbf{Doctoral
Advisor:} Nikolai Makarov (Caltech; undergraduate advisor: Viktor Havin)
\textbf{Fields Medal Citation:} \emph{``For the proof of conformal
invariance of percolation and the planar Ising model in statistical
physics.''}

\begin{center}\rule{0.5\linewidth}{0.5pt}\end{center}

\section{Table of Contents}\label{ch2010stanislavsmirnov:table-of-contents}

\begin{enumerate}
\def\labelenumi{\arabic{enumi}.}
\tightlist
\item
 Introduction
\item
 Biographical Background and Education
\item
 Statistical Mechanics: Background and Context
\item
 Percolation Theory
\item
 Conformal Invariance in Statistical Mechanics
\item
 Schramm--Loewner Evolution (SLE)
\item
 Smirnov's Proof of Conformal Invariance for Percolation
\item
 Critical Percolation on the Triangular Lattice
\item
 The Cardy--Smirnov Formula
\item
 The Planar Ising Model
\item
 Conformal Invariance of the Ising Model
\item
 The Discrete Complex Analysis Approach
\item
 Universality
\item
 Random Mats and the Gaussian Free Field
\item
 Further Results in Statistical Mechanics
\item
 Broader Impact and Influence
\item
 Awards and Recognition
\item
 Selected Major Publications
\item
 Ongoing Work and Future Directions
\item
 Conclusion
\item
 References and Further Reading
\end{enumerate}

\begin{center}\rule{0.5\linewidth}{0.5pt}\end{center}

\section{1. Introduction}\label{ch2010stanislavsmirnov:introduction}

Stanislav Smirnov was awarded the Fields Medal in 2010 for his
groundbreaking proof of conformal invariance for critical percolation
and the planar Ising model --- two of the most fundamental models in
statistical mechanics. His work resolved long-standing conjectures in
mathematical physics and provided a rigorous foundation for the
connection between discrete statistical models and continuous conformal
field theory.

The key idea in Smirnov's work is the use of \textbf{discrete complex
analysis} --- a general framework that allows one to study discrete
statistical models using complex-analytic methods. By showing that
certain discrete functions (observables in statistical models) converge
to holomorphic functions in the scaling limit, Smirnov was able to prove
conformal invariance and compute critical exponents exactly.

Smirnov's work is notable for its combination of mathematical rigor and
physical insight. The models he studied --- percolation and the Ising
model --- are among the important in statistical mechanics, and
their conformal invariance had been predicted by physicists for decades.
Smirnov provided the first rigorous proofs.

\begin{center}\rule{0.5\linewidth}{0.5pt}\end{center}

\section{2. Biographical Background and
Education}\label{ch2010stanislavsmirnov:biographical-background-and-education}

\subsection{2.1 Early Life}\label{ch2010stanislavsmirnov:early-life}

Stanislav Smirnov was born on September 3, 1970, in Leningrad (now
St.~Petersburg), USSR. He showed early mathematical talent and attended
the famous 239th Physics and Mathematics School in Leningrad, one of the
leading specialized schools in the Soviet Union.

\subsection{2.2 Mathematical Olympiads}\label{ch2010stanislavsmirnov:mathematical-olympiads}

Smirnov excelled in mathematical competitions. He won gold medals at the
International Mathematical Olympiad in 1986 (Warsaw) and 1987 (Havana)
--- one of only a handful of mathematicians to win two IMO gold medals.

\subsection{2.3 University Education}\label{ch2010stanislavsmirnov:university-education}

Smirnov entered Leningrad State University (now St.~Petersburg State
University) in 1987, where he studied mathematics. After the dissolution
of the Soviet Union, he continued his studies at the California
Institute of Technology (1990--1992) and then at Cornell University.

\subsection{2.4 Doctoral Studies}\label{ch2010stanislavsmirnov:doctoral-studies}

Smirnov earned his PhD in 1996 from the California Institute of
Technology under the supervision of Nikolai Makarov (having previously
studied under Viktor Havin at St.~Petersburg State University). His PhD
thesis, ``Spectral Analysis of Julia Sets,'' dealt with complex dynamics
--- a significant departure from his later work in statistical
mechanics.

However, the mathematical seeds of his later work were already present:
he was closely interested in complex analysis and its applications to
other areas of mathematics.

\subsection{2.5 Academic Career}\label{ch2010stanislavsmirnov:academic-career}

After his PhD, Smirnov held positions at: - \textbf{The Royal Institute
of Technology (KTH), Stockholm:} Postdoctoral researcher (1997--1998) -
\textbf{Max Planck Institute for Mathematics, Bonn:} Postdoctoral
researcher (1998--1999) - \textbf{Yale University:} Assistant Professor
(1999--2001) - \textbf{University of Geneva:} Associate Professor
(2001--2003), Professor (2003--present) - \textbf{St.~Petersburg State
University:} Head of the Laboratory of Mathematical Physics, Chebyshev
Laboratory (2008--present)

He maintains dual appointments in Geneva and St.~Petersburg.

\begin{center}\rule{0.5\linewidth}{0.5pt}\end{center}

\section{3. Statistical Mechanics: Background and
Context}\label{ch2010stanislavsmirnov:statistical-mechanics-background-and-context}

\subsection{3.1 The Goal of Statistical
Mechanics}\label{ch2010stanislavsmirnov:the-goal-of-statistical-mechanics}

Statistical mechanics aims to understand macroscopic properties of
matter from microscopic interactions. For a system of many interacting
particles on a lattice, one studies the \textbf{partition function}:

\[
Z = \sum_{\text{configurations}} e^{-\beta E(\text{configuration})}
\]

where \(\beta = 1/(k_B T)\) is the inverse temperature.

\subsection{3.2 Phase Transitions and Critical
Phenomena}\label{ch2010stanislavsmirnov:phase-transitions-and-critical-phenomena}

At certain critical temperatures, statistical systems undergo phase
transitions, characterized by: - \textbf{Divergent correlation lengths:}
Correlations decay as power laws rather than exponentially. -
\textbf{Scale invariance:} The system looks the same at all length
scales. - \textbf{Universality:} Different microscopic models exhibit
the same critical behavior.

\subsection{3.3 Conformal Invariance}\label{ch2010stanislavsmirnov:conformal-invariance}

At the critical point, scale invariance is typically enhanced to
\textbf{conformal invariance} --- invariance under angle-preserving
transformations. This was predicted by Belavin, Polyakov, and
Zamolodchikov in 1984 through conformal field theory (CFT), which
provides exact predictions for critical exponents and correlation
functions in two dimensions.

Proving conformal invariance from microscopic models was one of the
central challenges of mathematical physics.

\begin{center}\rule{0.5\linewidth}{0.5pt}\end{center}

\section{4. Percolation Theory}\label{ch2010stanislavsmirnov:percolation-theory}

\subsection{4.1 What Is Percolation?}\label{ch2010stanislavsmirnov:what-is-percolation}

Percolation is the simplest model of a disordered medium. Consider a
graph (e.g., the triangular lattice), where each edge is independently
open with probability \(p\) and closed with probability \(1-p\). The
question is: does there exist an infinite connected cluster of open
edges?

The \textbf{percolation threshold} \(p_c\) is the critical probability
above which an infinite cluster appears almost surely.

\subsection{4.2 Critical Percolation}\label{ch2010stanislavsmirnov:critical-percolation}

At \(p = p_c\), the system is at its critical point. For two-dimensional
percolation on the triangular lattice, \(p_c = 1/2\). At criticality,
the system exhibits: - \textbf{Fractal clusters:} Cluster boundaries are
fractal, with Hausdorff dimension 7/4. - \textbf{Power-law
correlations:} The probability that two distant points are connected
decays as a power law. - \textbf{Conformal invariance:} Predicted by
Cardy (1992) and others.

\subsection{4.3 Physical Importance}\label{ch2010stanislavsmirnov:physical-importance}

Percolation is the simplest model exhibiting a phase transition, and it
appears in numerous physical contexts: - Flow in porous media - Spread
of disease - Forest fires - Conductivity in random networks - Social
networks

\begin{center}\rule{0.5\linewidth}{0.5pt}\end{center}

\section{5. Conformal Invariance in Statistical
Mechanics}\label{ch2010stanislavsmirnov:conformal-invariance-in-statistical-mechanics}

\subsection{5.1 The Prediction}\label{ch2010stanislavsmirnov:the-prediction}

In 1992, John Cardy used conformal field theory to predict the exact
formula for the probability that a cluster connects two disjoint arcs on
the boundary of a simply connected domain --- the \textbf{Cardy
formula}. In a domain \(\Omega\) with four marked boundary points
\(a, b, c, d\), the probability that percolation connects the arc
\(\widehat{ab}\) to \(\widehat{cd}\) is:

\[
P(a, b, c, d) = \frac{\Gamma(2/3)}{\Gamma(1/3)^2} \cdot \eta^{1/3} \cdot {}_2F_1(1/3, 2/3; 4/3; \eta)
\]

where \(\eta\) is the cross-ratio of the four points.

\subsection{5.2 The Challenge}\label{ch2010stanislavsmirnov:the-challenge}

Cardy's formula was derived using non-rigorous conformal field theory
methods. Proving it from the microscopic definition of percolation
required: 1. Showing that crossing probabilities have a scaling limit 2.
Showing that this limit is conformally invariant 3. Computing the limit
explicitly

This seemed extremely difficult, as discrete percolation has no obvious
conformal symmetry.

\subsection{5.3 Smirnov's Approach}\label{ch2010stanislavsmirnov:smirnovs-approach}

Smirnov introduced a new method: \textbf{discrete complex analysis}. He
defined a \textbf{discrete holomorphic observable} --- a function on the
lattice that is approximately holomorphic --- and showed that in the
scaling limit, it converges to a holomorphic function satisfying certain
boundary conditions. This function turns out to be a conformal map, and
its values determine crossing probabilities.

\begin{center}\rule{0.5\linewidth}{0.5pt}\end{center}

\section{6. Schramm--Loewner Evolution
(SLE)}\label{ch2010stanislavsmirnov:schrammloewner-evolution-sle}

\subsection{6.1 Introduction}\label{ch2010stanislavsmirnov:introduction-1}

Schramm--Loewner Evolution, introduced by Oded Schramm in 2000, is a
family of random curves in the plane that are conformally invariant and
satisfy a certain Markov property. SLE is parameterized by a real number
\(\kappa > 0\), and different values of \(\kappa\) correspond to
different statistical mechanics models: - \(\kappa = 2\): Loop-erased
random walk - \(\kappa = 8/3\): Percolation cluster boundaries -
\(\kappa = 3\): Ising model interfaces - \(\kappa = 4\): Gaussian free
field level lines - \(\kappa = 6\): Percolation exploration paths -
\(\kappa = 8\): Uniform spanning tree

\subsection{6.2 SLE and Conformal
Invariance}\label{ch2010stanislavsmirnov:sle-and-conformal-invariance}

SLE provides a way to classify and study conformally invariant random
curves in the plane. If one can prove that the scaling limit of a
discrete model is SLE for an appropriate \(\kappa\), then many
properties of the model can be computed exactly.

\subsection{6.3 Smirnov's Results on SLE}\label{ch2010stanislavsmirnov:smirnovs-results-on-sle}

Smirnov's proof of conformal invariance for percolation implies that the
scaling limit of percolation exploration paths is SLE\(_6\). For the
Ising model, the scaling limit is SLE\(_3\) (for the spin interface) and
SLE\(_{16/3}\) (for the FK cluster boundaries).

\begin{center}\rule{0.5\linewidth}{0.5pt}\end{center}

\section{7. Smirnov's Proof of Conformal Invariance for
Percolation}\label{ch2010stanislavsmirnov:smirnovs-proof-of-conformal-invariance-for-percolation}

\subsection{7.1 Discrete Holomorphic
Functions}\label{ch2010stanislavsmirnov:discrete-holomorphic-functions}

A function \(f\) on the triangular lattice is \textbf{discrete
holomorphic} if it satisfies a discrete version of the Cauchy--Riemann
equations. On the triangular lattice (which consists of equilateral
triangles), a natural definition is:

For a vertex \(v\) with neighbors \(v_1, v_2, v_3\) in counterclockwise
order:

\[
f(v_1)(\overline{v_2} - \overline{v_3}) + f(v_2)(\overline{v_3} - \overline{v_1}) + f(v_3)(\overline{v_1} - \overline{v_2}) = 0
\]

This is the discrete analogue of the contour integral around a small
loop being zero.

\subsection{7.2 The Percolation
Observable}\label{ch2010stanislavsmirnov:the-percolation-observable}

Smirnov defined a \textbf{local observable} \(H(z)\) for percolation on
a triangular lattice domain. For a site \(z\) on the lattice, \(H(z)\)
is (roughly) the probability that a certain percolation configuration
connects \(z\) to the boundary in a particular way.

\subsection{7.3 Holomorphicity}\label{ch2010stanislavsmirnov:holomorphicity}

Smirnov proved that \(H(z)\) is discrete holomorphic on the triangular
lattice. The proof uses the following key properties of site percolation
on the triangular lattice: - \textbf{Self-duality:} The model is
invariant under the transformation \(p \leftrightarrow 1-p\), and at
\(p = 1/2\), the model is self-dual. - \textbf{Locality:} The observable
satisfies certain local relations arising from the percolation
configuration.

\subsection{7.4 Convergence}\label{ch2010stanislavsmirnov:convergence}

Having established that \(H(z)\) is discrete holomorphic, Smirnov used
standard results about discrete holomorphic functions to show that, as
the lattice spacing goes to zero, \(H(z)\) converges to a holomorphic
function \(h(z)\) on the domain. The boundary conditions determine
\(h(z)\) uniquely: it is the conformal map from the domain to an
equilateral triangle.

\subsection{7.5 The Cardy Formula}\label{ch2010stanislavsmirnov:the-cardy-formula}

The limit function \(h(z)\) gives crossing probabilities. Specifically,
for a domain \(\Omega\) with four marked boundary points, the
probability that percolation connects the two arcs is given by the ratio
of sides of the equilateral triangle, which yields the Cardy formula.

\begin{center}\rule{0.5\linewidth}{0.5pt}\end{center}

\section{8. Critical Percolation on the Triangular
Lattice}\label{ch2010stanislavsmirnov:critical-percolation-on-the-triangular-lattice}

\subsection{8.1 Why the Triangular
Lattice?}\label{ch2010stanislavsmirnov:why-the-triangular-lattice}

Smirnov's proof relies on the special structure of the triangular
lattice: - It is the only regular lattice where site percolation is
exactly solvable at \(p_c = 1/2\). - The lattice admits a natural
discrete complex structure (the faces are equilateral triangles, so
rotations by \(60^\circ\) preserve the lattice). - Self-duality holds
exactly.

\subsection{8.2 Extensions}\label{ch2010stanislavsmirnov:extensions}

While Smirnov's original proof worked specifically for site percolation
on the triangular lattice, subsequent work has extended conformal
invariance to: - Bond percolation on the square lattice (using different
methods) - Percolation with different boundary conditions - Other
two-dimensional lattices (via universality arguments)

\subsection{\texorpdfstring{8.3 The Exceptional Nature of
\(p = 1/2\)}{8.3 The Exceptional Nature of p = 1/2}}\label{ch2010stanislavsmirnov:the-exceptional-nature-of-p-12}

Smirnov's proof crucially uses \(p = 1/2\). For \(p \neq 1/2\), the
model is not critical and does not exhibit conformal invariance. The
proof shows that at the critical point, the discrete holomorphic
observable is exactly holomorphic, not just approximately.

\begin{center}\rule{0.5\linewidth}{0.5pt}\end{center}

\section{9. The Cardy--Smirnov Formula}\label{ch2010stanislavsmirnov:the-cardysmirnov-formula}

\subsection{9.1 Precise Statement}\label{ch2010stanislavsmirnov:precise-statement}

The Cardy--Smirnov formula gives the probability that there is a
percolation crossing between two disjoint boundary arcs in a simply
connected domain. For a domain \(\Omega\) with boundary points
\(a, b, c, d\) in counterclockwise order:

\[
\mathbb{P}(\text{arc } ab \text{ connected to arc } cd) = \frac{\Gamma(2/3)}{\Gamma(1/3)^2} \cdot \eta^{1/3} \cdot {}_2F_1(1/3, 2/3; 4/3; \eta)
\]

where \(\eta\) is the conformal modulus (cross-ratio) of the four
points.

\subsection{9.2 Explicit Form for Specific
Domains}\label{ch2010stanislavsmirnov:explicit-form-for-specific-domains}

For a rectangle \([0, L] \times [0, 1]\), the probability that there is
a left-to-right percolation crossing is:

\[
\mathbb{P}(\text{crossing}) = \frac{\Gamma(2/3)}{\Gamma(1/3)^2} \cdot \left(\frac{L}{1+L}\right)^{1/3} \cdot {}_2F_1\left(\frac{1}{3}, \frac{2}{3}; \frac{4}{3}; \frac{L}{1+L}\right)
\]

\subsection{9.3 Numerical Verification}\label{ch2010stanislavsmirnov:numerical-verification}

The Cardy--Smirnov formula has been extensively verified numerically,
and the agreement between theory and simulation is excellent. It
provides one of the most precise tests of conformal invariance in
statistical mechanics.

\begin{center}\rule{0.5\linewidth}{0.5pt}\end{center}

\section{10. The Planar Ising Model}\label{ch2010stanislavsmirnov:the-planar-ising-model}

\subsection{10.1 The Model}\label{ch2010stanislavsmirnov:the-model}

The Ising model is a model of ferromagnetism on a lattice. Each site
\(i\) carries a spin \(\sigma_i \in \{\pm 1\}\). The energy of a
configuration is:

\[
E(\sigma) = -J \sum_{\langle i,j \rangle} \sigma_i \sigma_j
\]

where the sum is over nearest neighbors. The model undergoes a
ferromagnetic phase transition at the \textbf{Curie temperature}
\(T_c\).

\subsection{10.2 Exact Solution}\label{ch2010stanislavsmirnov:exact-solution}

The two-dimensional Ising model was solved exactly by Lars Onsager in
1944, who computed the spontaneous magnetization and critical exponents.
This was one of the greatest achievements of 20th-century mathematical
physics.

\subsection{10.3 Conformal Invariance of the Ising
Model}\label{ch2010stanislavsmirnov:conformal-invariance-of-the-ising-model}

At the critical temperature, the Ising model is predicted to be
conformally invariant. This was the basis for the famous BPZ paper
(Belavin, Polyakov, Zamolodchikov, 1984) that founded conformal field
theory.

However, proving conformal invariance from the lattice model remained
open for over 60 years after Onsager's solution.

\begin{center}\rule{0.5\linewidth}{0.5pt}\end{center}

\section{11. Conformal Invariance of the Ising
Model}\label{ch2010stanislavsmirnov:conformal-invariance-of-the-ising-model-1}

\subsection{11.1 Smirnov's Proof}\label{ch2010stanislavsmirnov:smirnovs-proof}

In a series of papers (2005--2010), Smirnov proved conformal invariance
of the critical Ising model. The proof follows a similar strategy to his
percolation proof:

\begin{enumerate}
\def\labelenumi{\arabic{enumi}.}
\item
 \textbf{Define a discrete observable:} For the Ising model, the
 observable is related to the \textbf{order-disorder correlation
 function} --- a quantity introduced by Kadanoff and Ceva in 1971.
\item
 \textbf{Prove discrete holomorphicity:} The observable satisfies
 discrete Cauchy--Riemann equations at criticality. This follows from
 the \textbf{Kramers--Wannier duality} of the Ising model.
\item
 \textbf{Take the scaling limit:} As the lattice spacing goes to zero,
 the observable converges to a holomorphic function.
\item
 \textbf{Identify the limit:} The limit function is determined by the
 geometry of the domain, and its properties give exact formulas for
 correlation functions and interface curves.
\end{enumerate}

\subsection{11.2 The Ising Observable}\label{ch2010stanislavsmirnov:the-ising-observable}

For the Ising model on a domain \(\Omega\) with boundary conditions,
Smirnov defined:

\[
F(z) = \langle \sigma(z) \mu(\partial \Omega) \rangle
\]

where \(\sigma(z)\) is the spin operator at \(z\) and
\(\mu(\partial \Omega)\) is a disorder operator on the boundary. This
correlation function satisfies a discrete Cauchy--Riemann equation.

\subsection{11.3 Results}\label{ch2010stanislavsmirnov:results}

Smirnov's proof established: - \textbf{Conformal invariance of spin
correlations:} The scaling limit of spin correlations is conformally
invariant. - \textbf{SLE\(_3\) for spin interfaces:} The interfaces
between \(+\) and \(-\) spins converge to SLE\(_3\). -
\textbf{SLE\(_{16/3}\) for FK cluster boundaries:} The boundaries of
Fortuin--Kasteleyn clusters converge to SLE\(_{16/3}\).

\subsection{11.4 Complete Proof}\label{ch2010stanislavsmirnov:complete-proof}

The complete proof of conformal invariance for the Ising model was a
collaborative effort. Smirnov proved the key convergence results for the
discrete observable, and subsequent work by Chelkak, Hongler, Izyurov,
and others completed the program by establishing convergence of
correlations and interface curves.

\begin{center}\rule{0.5\linewidth}{0.5pt}\end{center}

\section{12. The Discrete Complex Analysis
Approach}\label{ch2010stanislavsmirnov:the-discrete-complex-analysis-approach}

\subsection{12.1 Discrete Holomorphic
Functions}\label{ch2010stanislavsmirnov:discrete-holomorphic-functions-1}

Smirnov's work established discrete complex analysis as a fundamental
tool in probability and statistical mechanics. A function \(f\) on a
planar lattice is discrete holomorphic if it satisfies:

\[
\oint_{\partial \triangle} f(z) \, dz = 0
\]

for every elementary triangle of the lattice, where the integral is
replaced by a sum over edges.

\subsection{12.2 Properties}\label{ch2010stanislavsmirnov:properties}

Discrete holomorphic functions share many properties with their
continuous counterparts: - \textbf{Maximum principle:} The maximum of
\(|f|\) occurs on the boundary. - \textbf{Cauchy integral formula:}
Values inside a region can be recovered from boundary values. -
\textbf{Convergence:} Under appropriate conditions, sequences of
discrete holomorphic functions converge to holomorphic functions.

\subsection{12.3 Applications Beyond Percolation and
Ising}\label{ch2010stanislavsmirnov:applications-beyond-percolation-and-ising}

The discrete complex analysis approach has been applied to many other
models: - \textbf{Random cluster model:} Conformal invariance for the FK
model - \textbf{Gaussian free field:} Discrete harmonic functions
converge to the continuum GFF - \textbf{Loop-erased random walk:}
Convergence to SLE\(_2\) (Lawler, Schramm, Werner) - \textbf{Height
functions:} Convergence to Gaussian free field

\begin{center}\rule{0.5\linewidth}{0.5pt}\end{center}

\section{13. Universality}\label{ch2010stanislavsmirnov:universality}

\subsection{13.1 The Universality
Hypothesis}\label{ch2010stanislavsmirnov:the-universality-hypothesis}

Universality is the idea that many different microscopic models exhibit
the same critical behavior. For example, percolation on the triangular
lattice and percolation on the square lattice should have the same
critical exponents.

\subsection{13.2 Smirnov's Contribution}\label{ch2010stanislavsmirnov:smirnovs-contribution}

Smirnov's proof of conformal invariance for specific models (triangular
lattice percolation, planar Ising) provides evidence for universality: -
By proving conformal invariance for one lattice, one can compute
critical exponents exactly. - These exponents are expected to be
universal --- the same for all lattices in the same universality class.
- Numerical simulations and heuristic arguments support this
expectation.

\subsection{13.3 Proving Universality}\label{ch2010stanislavsmirnov:proving-universality}

Rigorous proof of universality --- showing that critical exponents are
independent of the lattice --- remains a major open problem. Smirnov's
work provides the foundation for such proofs in two dimensions by
establishing conformal invariance for specific models.

\begin{center}\rule{0.5\linewidth}{0.5pt}\end{center}

\section{14. Random Mats and the Gaussian Free
Field}\label{ch2010stanislavsmirnov:random-mats-and-the-gaussian-free-field}

\subsection{14.1 Random Surfaces}\label{ch2010stanislavsmirnov:random-surfaces}

Smirnov has also worked on random surfaces and the Gaussian free field
(GFF). The GFF is the canonical random surface in two dimensions and
appears as the scaling limit of many discrete models.

\subsection{14.2 The Gaussian Free Field}\label{ch2010stanislavsmirnov:the-gaussian-free-field}

The Gaussian free field is a Gaussian random distribution on a domain
\(\Omega\) whose covariance is the Green's function of the Laplacian. It
plays a central role in: - Conformal field theory (it corresponds to the
massless free field) - Random surfaces - SLE theory (level lines of the
GFF are SLE\(_4\))

\subsection{14.3 Discrete Gaussian Free
Field}\label{ch2010stanislavsmirnov:discrete-gaussian-free-field}

Smirnov has studied the convergence of discrete height functions to the
GFF, using his discrete harmonic analysis techniques. This has
applications to: - The dimer model - The Gaussian free field on random
planar maps - Interfaces in the Ising model

\begin{center}\rule{0.5\linewidth}{0.5pt}\end{center}

\section{15. Further Results in Statistical
Mechanics}\label{ch2010stanislavsmirnov:further-results-in-statistical-mechanics}

\subsection{15.1 The Random Cluster
Model}\label{ch2010stanislavsmirnov:the-random-cluster-model}

The random cluster model (also known as the FK model) generalizes both
percolation and the Ising model. It is parameterized by \(p\) (edge
weight) and \(q\) (cluster weight): - \(q = 1\): Percolation -
\(q = 2\): Ising model (via the Fortuin--Kasteleyn representation)

Smirnov and collaborators have proved conformal invariance for the FK
model at criticality, extending his earlier results.

\subsection{\texorpdfstring{15.2 The \(O(n)\)
Model}{15.2 The O(n) Model}}\label{ch2010stanislavsmirnov:the-on-model}

The \(O(n)\) model is a family of loop models that includes percolation
(\(n = 1\)), the Ising model (\(n = 2\)), and self-avoiding walks
(\(n = 0\)). Smirnov's methods have been applied to certain values of
\(n\), with further progress expected.

\subsection{15.3 The Dimer Model}\label{ch2010stanislavsmirnov:the-dimer-model}

The dimer model (perfect matchings on a lattice) is exactly solvable and
exhibits conformal invariance. Smirnov's discrete complex analysis
approach provides a rigorous framework for understanding the dimer
model's scaling limit.

\begin{center}\rule{0.5\linewidth}{0.5pt}\end{center}

\section{16. Broader Impact and
Influence}\label{ch2010stanislavsmirnov:broader-impact-and-influence}

\subsection{16.1 Creation of a New Field}\label{ch2010stanislavsmirnov:creation-of-a-new-field}

Smirnov's work, together with the development of SLE by Schramm, Lawler,
and Werner, created the modern field of \textbf{random conformal
geometry}. This field studies the scaling limits of discrete statistical
mechanics models using the tools of complex analysis and probability.

\subsection{16.2 Influence on Young
Researchers}\label{ch2010stanislavsmirnov:influence-on-young-researchers}

Smirnov has supervised numerous PhD students and postdocs, many of whom
have become leaders in the field: - \textbf{Dmitry Chelkak:} Extensions
of discrete complex analysis and the Ising model - \textbf{Konstantin
Izyurov:} Conformal invariance of the Ising model - \textbf{Hugo
Duminil-Copin:} Percolation and phase transitions (Fields Medal 2022)

\subsection{16.3 The Geneva School}\label{ch2010stanislavsmirnov:the-geneva-school}

The University of Geneva, where Smirnov is based, has become a major
center for random conformal geometry. Smirnov has built a strong group
there, including: - Hugo Duminil-Copin (until his move to IHÉS) - Dmitry
Chelkak - Numerical and computational collaborators

\subsection{16.4 The Smirnov--Duminil-Copin
Connection}\label{ch2010stanislavsmirnov:the-smirnovduminil-copin-connection}

Hugo Duminil-Copin, who won the Fields Medal in 2022, was a PhD student
of Smirnov at the University of Geneva. Duminil-Copin's work on phase
transitions in statistical mechanics builds directly on Smirnov's
methods.

\begin{center}\rule{0.5\linewidth}{0.5pt}\end{center}

\section{17. Awards and Recognition}\label{ch2010stanislavsmirnov:awards-and-recognition}

\textbf{2010:} - Fields Medal (International Mathematical Union)

\textbf{2009:} - European Mathematical Society Prize

\textbf{2006:} - Clay Research Award (with D. Chelkak) - Prize of the
St.~Petersburg Mathematical Society

\textbf{2004:} - Rollo Davidson Prize (University of Cambridge)

\textbf{2001:} - Prize of the Swiss Mathematical Society

\textbf{1998:} - Delwart Prize (Caltech)

\textbf{Honorary Positions:} - Corresponding Member, Russian Academy of
Sciences - Member, Academia Europaea

\begin{center}\rule{0.5\linewidth}{0.5pt}\end{center}

\section{18. Selected Major
Publications}\label{ch2010stanislavsmirnov:selected-major-publications}

\begin{enumerate}
\def\labelenumi{\arabic{enumi}.}
\item
 \textbf{``Critical percolation in the plane.''} arXiv:0909.4499, 2009.
 --- \emph{Comprehensive proof of conformal invariance for
 percolation.}
\item
 \textbf{``Conformal invariance of the critical Ising model.''}
 \emph{Proceedings of the International Congress of Mathematicians},
 2006. --- \emph{Overview of the Ising model result.}
\item
 \textbf{``Discrete complex analysis and probability.''}
 \emph{Proceedings of the International Congress of Mathematicians},
 2010. --- \emph{Survey of discrete complex analysis methods.}
\item
 \textbf{``Conformal invariance in random cluster models. I.
 Holomorphic fermions in the Ising model''} (with D. Chelkak).
 \emph{Annals of Mathematics} 175(3), 1--48, 2012. --- \emph{Detailed
 Ising model proof.}
\item
 \textbf{``Universality in the 2D Ising model and conformal invariance
 of fermionic observables''} (with D. Chelkak and H. Duminil-Copin).
 \emph{Inventiones Mathematicae} 189(3), 515--580, 2012. ---
 \emph{Universality for the Ising model.}
\item
 \textbf{``The Ising model on a random planar lattice: a discrete
 approach.''} \emph{Journal of Statistical Physics} 124(2--4),
 465--483, 2006.
\item
 \textbf{``Limit functions for the Ising model''} (with D. Chelkak).
 \emph{Duke Mathematical Journal} 162(1), 79--125, 2013.
\item
 \textbf{``Percolation crossing probabilities and conformal
 invariance.''} \emph{Comptes Rendus Mathématique} 333(3), 239--244,
 2001. --- \emph{Announcement of the percolation result.}
\end{enumerate}

\begin{center}\rule{0.5\linewidth}{0.5pt}\end{center}

\section{19. Ongoing Work and Future
Directions}\label{ch2010stanislavsmirnov:ongoing-work-and-future-directions}

\subsection{19.1 Complete Universality}\label{ch2010stanislavsmirnov:complete-universality}

Proving universality for percolation (that all two-dimensional lattices
have the same critical exponents) remains an open problem. Smirnov's
proof for the triangular lattice is the best evidence, but a general
proof is lacking.

\subsection{19.2 Self-Avoiding Walks}\label{ch2010stanislavsmirnov:self-avoiding-walks}

The scaling limit of self-avoiding walks (SAWs) is predicted to be
SLE\(_{8/3}\). Proving this remains one of the major open problems in
random conformal geometry.

\subsection{19.3 Higher-Dimensional
Models}\label{ch2010stanislavsmirnov:higher-dimensional-models}

While Smirnov's methods are inherently two-dimensional, there is ongoing
work to extend some ideas to higher dimensions. However, conformal
invariance is special to two dimensions, so different methods are
needed.

\subsection{19.4 Random Planar Maps}\label{ch2010stanislavsmirnov:random-planar-maps}

Smirnov has been working on random planar maps and their connection to
Liouville quantum gravity --- a direction that unifies random geometry,
conformal field theory, and probability.

\begin{center}\rule{0.5\linewidth}{0.5pt}\end{center}

\section{20. Conclusion}\label{ch2010stanislavsmirnov:conclusion}

Stanislav Smirnov's proof of conformal invariance for critical
percolation and the planar Ising model represents one of the greatest
achievements in 21st-century mathematical physics. His work resolved
conjectures that had stood for decades and provided the first rigorous
foundation for the connection between discrete statistical models and
conformal field theory.

The key to Smirnov's success was the development of \textbf{discrete
complex analysis} --- a framework that allows one to study discrete
models using the tools of complex analysis. This framework has since
been applied to numerous other models and has become a standard tool in
probability and statistical mechanics.

The Fields Medal citation --- \emph{``for the proof of conformal
invariance of percolation and the planar Ising model in statistical
physics''} --- captures the essence of his contribution. Smirnov did not
merely prove two important conjectures; he introduced methods that
transformed the field and opened up new directions of research.

His influence extends through his students, especially Hugo
Duminil-Copin (Fields Medal 2022), who has continued and extended the
Geneva tradition of rigorous statistical mechanics. The methods Smirnov
developed will continue to shape the field for decades to come.

\begin{center}\rule{0.5\linewidth}{0.5pt}\end{center}

\section{21. References and Further
Reading}\label{ch2010stanislavsmirnov:references-and-further-reading}

\subsection{Official Sources}\label{ch2010stanislavsmirnov:official-sources}

\begin{itemize}
\tightlist
\item
 ICM 2010 Fields Medal Announcement:
 https://www.icm2010.in/prize-winners-2010/fields-medal-stanislav-smirnov
\item
 IMU Fields Medal Citations:
 https://www.mathunion.org/imu-awards/fields-medal/fields-medal-2010
\end{itemize}

\subsection{Biographical and Overview
Articles}\label{ch2010stanislavsmirnov:biographical-and-overview-articles}

\begin{itemize}
\tightlist
\item
 Smirnov, S. ``Discrete Complex Analysis and Probability.'' ICM 2010
 Plenary Lecture.
\item
 Lawler, G. ``The Work of Stanislav Smirnov.'' ICM 2010 Laudatio.
\item
 Bauer, M. and Bernard, D. ``2D Growth Processes: SLE and Loewner
 Chains.'' \emph{Physics Reports} 432(3), 2006.
\end{itemize}

\subsection{Key Technical References}\label{ch2010stanislavsmirnov:key-technical-references}

\begin{itemize}
\tightlist
\item
 Smirnov, S. ``Critical percolation in the plane.'' arXiv:0909.4499.
\item
 Chelkak, D. and Smirnov, S. ``Universality in the 2D Ising model and
 conformal invariance of fermionic observables.'' \emph{Inventiones
 Mathematicae} 189 (2012), 515--580.
\item
 Cardy, J. ``Critical percolation in finite geometries.'' \emph{Journal
 of Physics A} 25 (1992), L201--L206.
\item
 Schramm, O. ``Scaling limits of loop-erased random walks and uniform
 spanning trees.'' \emph{Israel Journal of Mathematics} 118 (2000),
 221--288.
\end{itemize}

\subsection{Background Reading}\label{ch2010stanislavsmirnov:background-reading}

\begin{itemize}
\tightlist
\item
 Kager, W. and Nienhuis, B. ``A Guide to Stochastic Löwner Evolution
 and Its Applications.'' \emph{Journal of Statistical Physics} 115
 (2004), 1149--1229.
\item
 Lawler, G. \emph{Conformally Invariant Processes in the Plane.}
 American Mathematical Society, 2005.
\item
 Werner, W. ``Random Planar Curves and Schramm--Loewner Evolutions.''
 \emph{Lectures on Probability Theory and Statistics}, Springer, 2004.
\end{itemize}
